\section{Complex Analysis}

Here, we don't include proofs.

\subsection{Basics of Holomorphic Functions}

\begin{theorem}
    Let $f \colon \Omega \rightarrow \C$, where $\Omega \subseteq \C$ is open.
    The following statements are equivalent:
    \begin{itemize}
        \item $f$ is holomorphic at $a \in \Omega$.
        \item $f$ has a convergent power series $\sum_{n \geq 0} c_n(z - a)^n$ in an open 
            disc centered at $a \in \Omega$.
    \end{itemize}
\end{theorem}

\begin{df}
    A complex-valued function is entire if it is holomorphic over $\C$.
\end{df}

\begin{df}
    Consider $f(z) = u + iv$ be a complex-valued function. Writing $z = x + iy$, the 
    \emph{Cauchy-Riemann equations} refer to the set of differential equations
    \[
    \pfrac{u}{x} = \pfrac{v}{y}, \qquad \pfrac{u}{y} = -\pfrac{v}{x}.
    \]
\end{df}

\begin{theorem}
    Consider $f = u + iv \colon \Omega \rightarrow \C$, defomed on an open set $\Omega \subseteq \C$.
    If $u$ and $v$ are continuously differentiable and satisfy the Cauchy-Riemann equations on $\Omega$,
    then $f$ is holomorphic on $\Omega$ and $f'(z) = \frac{\partial f(z)}{\partial z}$.
\end{theorem}

\begin{example}
    $f(z) = \bar{z}$ is not holomorphic on $\C$. $\bar{z} = x - iy$, 
    so $u_x' = 1 \neq -1 = v_y'$.
\end{example}

\begin{exercise}
Suppose that $f(x + iy) = u(x, y) + iv(x, y)$. If $u(x, y) = x^2 - y^2$ and $v(1,1) = 2$. Find $v(4,1)$.
\end{exercise}
\begin{solution}
The key is the FTC and the Cuachy-Riemann equations, because $f$ is clearly holomorphic.
\[
v(4,1) - v(1,1) = \int_1^4 v_x'(x,1) dx = \int_1^4 - u_y'(x,1) dy = \int_1^4 2dy = 6.
\]
Therefore, $v(4,1) = 8$. Note the use of the Cauchy Riemann equations.
\end{solution}

\subsection{Cauchy's Integral Formula}

\begin{theorem}[Cauchy]
    Let $f$ be holomorphic in an open set $\Omega$, then  $f$ has infinitely many 
    complex derivatives in $\Omega$. Moreover, for simple, closed, piecewise-smooth curve 
    $\gamma \subset \Omega$ and any $z$ lying inside $\gamma$ we have
    \[
    \frac{d^nf(z)}{dz^n} = \frac{n!}{2 \pi i} \oint_{\gamma} \frac{f(\eta)}{(\eta - z)^{n + 1}} d\eta.
    \]
\end{theorem}

\begin{exercise}
    Evaluate the following integral, for $\gamma = \{ z \colon |z| = 2\}$:
    \[
    \frac{1}{2 \pi i} \oint_{\gamma} \frac{e^z}{(z + i)(z - i)} dz
    \]
\end{exercise}
\begin{solution}
    We have as follows:
    \[
    \frac{1}{2 \pi i} \oint_{\gamma} \frac{e^z}{(z + i)(z - i)} dz = \frac{1}{2 \pi i} \frac{1}{2i} \oint_{\{ z \colon |z| = 2\}} \left( \frac{e^z}{z + i} - \frac{e^z}{z - i} \right) dz = \frac{1}{2i} (e^i - e^{-i}) = \sin 1.
    \]
\end{solution}

\begin{exercise}
Let $C = \{z = 1 + e^{i \theta} \colon \theta \in [0,2\pi]\}$. Calculate
\[
\frac{1}{2\pi i} \int_C \left(\frac{\sin z}{z - 1}\right)^2 dz.
\]
\end{exercise}
\begin{solution}
Use Cauchy's integral formula with $n = 1$. We want the first derivative of $\sin^2(z)$,
and plug in $1$ into it. We have $2 \sin(1)\cos(1) = \sin(2)$.
\end{solution}


\begin{theorem}[Morera]
    Suppose $f$ is a continuous function in the open disc $D$ such that 
    for any triangle $T$ contained in $D$
    \[
    \int_T f(z) dz = 0,
    \]
    then $f$ is holomorphic.
\end{theorem}


\begin{corollary}[Liouville]
    A bounded entire function is constant.
\end{corollary}

\subsection{Residue Theory}


\begin{df}
    Let $f(z) = \sum \limits_{-\infty}^{\infty} a_n(z - z_0)^n$, for $0 < |z - z_0| < R$
    be the Laurent series for $f$ at $z_0$. The \emph{residue} of $f$ at $z_0$ is $\Res [f, z_0] = a_{-1}$.
\end{df}
\iffalse
\begin{theorem}
    Let $\gamma \subset \{z \colon 0 < |z - z_0| < R \}$ be a simple, closed, 
    piecewise-smooth curve that contains $z_0$. Then,
    \[
    \Res [f, z_0] = \frac{1}{2 \pi i} \oint_{\gamma} f(z) dz.
    \]
\end{theorem}
\fi
\begin{theorem}
    Let $f$ be holomorphic function inside and on a simple, closed, 
    piecewise-smooth curve $\gamma$ except at the singularities $z_1, \dots, z_n$ in its interior. 
    Then,
    \[
    \oint_{\gamma} f(z) dz = 2 \pi i \sum \limits_{j = 1}^n \Res [f, z_j]
    \]
\end{theorem}

\begin{remark}
    Note that we can find the residue of $f$ at a pole $z_0$ of order $m$ as follows:
    \[
    \Res [f, z_0] = \lim \limits_{z \rightarrow z_0} \frac{1}{(m - 1)!} \frac{d^{m - 1}}{dz^{m - 1}} ((z - z_0)^{m}f(z)).
    \]
\end{remark}

\begin{exercise}
    Evaluate $\oint_{\gamma} \frac{1}{z^5 - z^3} dz$, where 
    $\gamma = \{ z \colon |z| = \frac{1}{2} \}$.
\end{exercise}
\begin{solution}
    Let $f(z) = \frac{1}{z^5 - z^3}$. The only relevant singularity in $\gamma$ is $z = 0$. 
    Thus, the integral is equal to $\Res[f,0]$, which is a pole of order $3$:
    \begin{align*}
    \lim_{z\rightarrow0} \frac{1}{2!}\frac{d^2}{dz^2} \left(z^3 \cdot f(z)\right)
    &= \frac{1}{2} \lim_{z \rightarrow 0} \left(\frac{-2z}{(z^2-1)^2}\right)' \\
    &= \frac{1}{2} \lim_{z \rightarrow 0} \frac{-2(z^2 - 1)^2 +4z(z^2 - 1)\cdot 2z}{(z^2-1)^4} \\\
    &= -1
    \end{align*}
    Thus, the integral evaluates to $-2\pi i$.
\end{solution}